\documentclass[aps,pra,onecolumn,superscriptaddress,nofootinbib,longbibliography]{revtex4-2}

\usepackage{amsmath}
\usepackage{amssymb}
\usepackage{amsthm}
\usepackage{booktabs}
\usepackage[colorlinks=true,linkcolor=blue,citecolor=blue,urlcolor=blue]{hyperref}

\theoremstyle{plain}
\newtheorem{theorem}{Theorem}
\newtheorem{lemma}[theorem]{Lemma}
\newtheorem{proposition}[theorem]{Proposition}
\newtheorem{corollary}[theorem]{Corollary}
\theoremstyle{definition}
\newtheorem{definition}[theorem]{Definition}
\theoremstyle{remark}
\newtheorem{remark}[theorem]{Remark}

\newcommand{\F}{\mathrm{F}}
\newcommand{\GHZ}{\mathrm{GHZ}_3}
\newcommand{\LC}{\mathrm{LC}}
\newcommand{\rw}{\mathrm{rw}}
\newcommand{\DHclass}{\mathsf{DH}}
\newcommand{\Prod}{\mathsf{Prod}}
\newcommand{\code}[1]{\texttt{#1}}
\newcommand{\nbr}{N}

\begin{document}

\title{Distance-hereditary graphs characterize minimum-fusion synthesis of graph states}

\author{Sean Howell}
\email{sean.g.howell@gmail.com}

\begin{abstract}
In fusion-based quantum computation a graph state $|G\rangle$ on $n$ qubits is assembled from
small entangled resource states by destructive two-qubit fusion measurements. We study the
minimum number $\F(G)$ of fusions needed when the resource is the three-qubit
Greenberger-Horne-Zeilinger state, every fusion is a Bell-type measurement on any two qubits,
and single-qubit Clifford operations are free. Counting qubits and connected components gives
$\F(G)\ge n-3$ for every connected graph. We prove that $\F(G)=n-3$ if and only if $G$ is
distance-hereditary, equivalently of rank-width at most one, equivalently free of
$C_5$-vertex-minors. The graph-theoretic core of the proof, that the graphs reachable at this
bound by pendant, twin and local-complementation steps are exactly the distance-hereditary
graphs, is a Lean 4 theorem. The steps connecting it to qubits are proved in the text: the
graph rule for a fusion of two disjoint graph states, the absorption of single-qubit Cliffords
on fused qubits into local complementations, and a normal form that reorders any schedule into
a single growing component. We also report a classification of all 585 connected
local-complementation orbits on at most nine vertices, certified by two independent stabilizer
engines, which agrees with the theorem on every orbit and gives nineteen exact values of
$\F$ above the bound, including a rank-width-two orbit that needs $n+3$ fusions, so rank-width
does not determine the cost above the extremal class. The characterization was proposed by a
language model working inside a harness that admits only machine-checked artifacts, from
labeled data on at most seven vertices, and the Lean proofs were written by the same model.
\end{abstract}

\maketitle

\section{Introduction}\label{sec:intro}

Fusion-based quantum computation (FBQC) builds fault-tolerant computation from a supply of
small, identical entangled \emph{resource states} that are joined by \emph{fusions}: destructive
joint measurements of two qubits~\cite{Bartolucci2023fbqc}. In a photonic architecture the
resource states and the fusions are the main physical costs, so for any target entangled state
a basic question is how few fusions suffice to make it. For graph states, the stabilizer states
in one-to-one correspondence with simple graphs~\cite{Hein2004multiparty,Hein2006graphstates},
this is a combinatorial optimization over graphs, and it has been studied numerically and with
dynamic programming in several resource models~\cite{Lobl2025generating,Lee2023graph,Ji2026motifs}.

We work in a specific model, which we call the $\GHZ$ model: the resource is the three-qubit GHZ
state, whose graph is the three-vertex path $P_3$; each fusion destructively measures a
commuting pair of two-qubit Pauli operators that distinguishes the four Bell states; and
single-qubit Clifford operations are free throughout. Let $\F(G)$ be the least number of fusions
that produce a state equal, up to single-qubit Cliffords, to the graph state $|G\rangle$. A short
count of qubits and connected components gives $\F(G)\ge n-3$ for a connected graph on $n\ge 3$
vertices (Proposition~\ref{prop:lb}). Our main result identifies the graphs that attain this
bound.

\begin{theorem}[Characterization]\label{thm:main}
For a connected graph $G$ on $n\ge 3$ vertices, in the $\GHZ$ model,
\[
  \F(G)=n-3 \quad\Longleftrightarrow\quad G \text{ is distance-hereditary.}
\]
\end{theorem}

Distance-hereditary graphs~\cite{Howorka1977characterization,Bandelt1986distance} are the graphs
in which every connected induced subgraph preserves distances. They are exactly the graphs of
rank-width at most one~\cite{Oum2005rankwidth}, the graphs with no vertex-minor isomorphic to
the five-cycle $C_5$~\cite{Bouchet1988transforming}, and the graphs obtained from a single vertex
by repeatedly adding pendant vertices and twins~\cite{Bandelt1986distance}. Rank-width is
invariant under local complementation~\cite{Oum2005rankwidth}, and $\F$ is invariant under
local Clifford operations, so Theorem~\ref{thm:main} matches the fusion cost at the bound to a
local-Clifford-invariant measure of the entanglement structure of $|G\rangle$.

\paragraph{What is proved, and how.} Because the result was obtained with machine assistance we
state plainly what kind of argument supports each part.
\begin{itemize}
\item \emph{Lean theorems.} The graph-theoretic core is checked by the Lean~4
  kernel~\cite{deMoura2021lean4} against \textsc{mathlib}~\cite{mathlib2020} with only the
  axioms \code{propext}, \code{Classical.choice} and \code{Quot.sound}: a graph is reachable from
  $P_3$ by $k$ pendant and false-twin steps together with free local complementations if and only
  if it is distance-hereditary with $k+3$ vertices (Theorem~\ref{thm:model}); the counting form
  of the lower bound; and $\F=n-3$ for paths, stars, all trees, complete graphs and complete
  bipartite graphs, which are special cases (Section~\ref{sec:lean}).
\item \emph{Written proofs.} The steps that connect the graph-level theorem to qubits are proved
  in the text: the graph rule for a fusion of two disjoint graph states
  (Lemma~\ref{lem:merge}); the absorption of single-qubit Cliffords on a fused qubit into local
  complementations (Lemma~\ref{lem:frame}); and a normal form that reorders any schedule into one
  that grows a single component (Lemma~\ref{lem:normal}). With these, Theorem~\ref{thm:main} has
  a complete proof (Section~\ref{sec:proof}).
\item \emph{Independent computation.} Two stabilizer engines that share no code agree with the
  fusion rule on random inputs, and an exhaustive search over all $585$ connected
  local-complementation orbits with $n\le 9$ agrees with Theorem~\ref{thm:main} on every orbit
  (Section~\ref{sec:table}).
\end{itemize}

\paragraph{How the result was found.} The characterization was not found by hand. A language
model, Fable~5, working inside a research harness that records nothing above the status of a
heuristic without a machine artifact and never lets the model execute code, classified the orbits,
proposed the distance-hereditary characterization from the $44$ labeled orbits with $n\le 7$, and
wrote the Lean proofs; the $541$ held-out orbits with $n\in\{8,9\}$ then agreed with the
proposal without exception (Section~\ref{sec:method}). We describe the harness briefly because
the provenance bears on how the results should be read, but the mathematics stands on its own:
every step of the proof is either a Lean theorem or written out here.

The paper is organized as follows. Section~\ref{sec:model} defines the model, proves the fusion
rule and the lower bound. Section~\ref{sec:proof} proves Theorem~\ref{thm:main}, with the proof
of the normal form in Appendix~\ref{app:normal}. Section~\ref{sec:lean} describes the Lean
development and Section~\ref{sec:table} the certified tablebase. Section~\ref{sec:method}
describes the harness and the discovery, Section~\ref{sec:related} relates the result to
previous work, and Section~\ref{sec:discussion} lists open questions.

\section{The model}\label{sec:model}

\paragraph{Graph states and local complementation.} For a simple graph $G=(V,E)$ the graph state
$|G\rangle=\prod_{\{u,v\}\in E}\mathrm{CZ}_{uv}\,|{+}\rangle^{\otimes V}$ is the unique common
$+1$ eigenvector of the operators $K_v=X_v\prod_{u\in\nbr(v)}Z_u$, $v\in V$, where $\nbr(v)$ is the
neighborhood of $v$. Local complementation at $v$ replaces the subgraph induced on $\nbr(v)$ by
its complement; we write $\tau_v(G)$. It is implemented on graph states by a local Clifford
unitary~\cite{VandenNest2004graphical},
\begin{equation}\label{eq:lc}
  |\tau_v(G)\rangle = e^{-i\pi X_v/4}\prod_{u\in\nbr(v)}e^{i\pi Z_u/4}\,|G\rangle,
\end{equation}
and two graph states on the same qubits are related by a product of single-qubit Clifford
unitaries if and only if their graphs are related by a sequence of local
complementations~\cite{VandenNest2004graphical}. The qubits that survive a schedule carry no fixed
labels, so we also identify isomorphic graphs, and write $G\sim_{\LC}H$ for the equivalence
generated by local complementation and isomorphism.

\paragraph{Resources, fusions and schedules.} A schedule starts from $g$ copies of the
three-qubit GHZ state; its graph is $P_3$, and it is $\LC$-equivalent to the triangle $K_3$. A
\emph{fusion} on two distinct qubits $a,b$ is the destructive measurement of the commuting
observables $X_aZ_b$ and $Z_aX_b$, after which $a$ and $b$ are discarded. This is a Bell
measurement up to a Hadamard on $b$; any other Bell-type convention differs from it by
single-qubit Cliffords on $a$ and $b$, which are free in the model. A schedule interleaves
fusions, on any pairs of qubits, with arbitrary single-qubit Clifford gates. It \emph{produces}
$G$ if the state of the surviving qubits is $\LC$-equivalent to $|G\rangle$. For a stabilizer
state the outcomes of a Pauli measurement give post-measurement states that differ by a Pauli
operator~\cite{Gottesman1997thesis}, so whether a schedule produces $G$ does not depend on the
outcomes; we treat fusions as deterministic, leaving the probabilistic success of linear-optical
fusion to a layer above this problem. The cost of a schedule is its number of fusions, and
$\F(G)$ is the least cost of a schedule that produces $G$.

\begin{lemma}[Fusion of disjoint graph states]\label{lem:merge}
Let $G_1$ and $G_2$ be graphs on disjoint vertex sets, $a$ a vertex of $G_1$ and $b$ a vertex of
$G_2$. A fusion on $(a,b)$ applied to $|G_1\rangle\otimes|G_2\rangle$ leaves, up to a Pauli
operator that depends on the outcomes, the graph state $|G_1\bowtie_{a,b}G_2\rangle$ whose vertex
set is $(V_1\cup V_2)\setminus\{a,b\}$ and whose edges are the edges of $G_1$ and $G_2$ not
incident to $a$ or $b$, together with every pair $\{v,w\}$ with $v\in\nbr_{G_1}(a)$ and
$w\in\nbr_{G_2}(b)$.
\end{lemma}

\begin{proof}
The state is stabilized by $K_v$ for all $v\in V_1\cup V_2$. The observable $X_aZ_b$
anticommutes exactly with $K_b$ and with $K_v$ for $v\in\nbr(a)$; the observable $Z_aX_b$
anticommutes exactly with $K_a$ and with $K_w$ for $w\in\nbr(b)$. Measuring $X_aZ_b$ with the
standard update rule~\cite{Gottesman1997thesis} replaces $K_v$ by $K_vK_b$ for $v\in\nbr(a)$ and
$K_b$ by $\pm X_aZ_b$. Each $K_vK_b$ commutes with $Z_aX_b$ (both factors contribute a
commuting component on $a$ and on $b$), so measuring $Z_aX_b$ next replaces $K_w$ by $K_wK_a$
for $w\in\nbr(b)$ and $K_a$ by $\pm Z_aX_b$. Multiplying $K_vK_b$ by $\pm Z_aX_b$ and $K_wK_a$ by
$\pm X_aZ_b$ removes their support on $a$ and $b$, leaving, up to sign,
\[
  X_v\,Z_{\nbr(v)\setminus\{a\}}\,Z_{\nbr(b)} \quad (v\in\nbr(a)),\qquad
  X_w\,Z_{\nbr(w)\setminus\{b\}}\,Z_{\nbr(a)} \quad (w\in\nbr(b)),
\]
and $K_u$ unchanged for every other surviving vertex $u$. These are the stabilizers $K_u$ of the
stated graph on the surviving qubits; the two measured observables stabilize a Bell-type state of
$a,b$, which is discarded. Signs are corrected by Pauli operators.
\end{proof}

Fusing a vertex $a$ of a graph with a qubit of a fresh resource therefore has one of three
effects, according to the graph used to represent the resource and the qubit fused (each is
available because single-qubit Cliffords on a fresh resource are free):
\begin{itemize}
\item a \emph{leaf} of $P_3$ (neighbor $c$, other leaf $\ell$): $a$ is replaced by $c$ with the same
  neighborhood, and $\ell$ is a new pendant vertex at $c$; up to relabeling, a \emph{pendant} is
  added at $a$;
\item the \emph{center} of $P_3$ (leaves $\ell_1,\ell_2$): $a$ is replaced by $\ell_1$ and $\ell_2$,
  each with the neighborhood of $a$ and not adjacent to each other; a \emph{false twin} is added
  at $a$;
\item a vertex of $K_3$ (other vertices $y,z$, adjacent): $a$ is replaced by the adjacent pair
  $y,z$; a \emph{true twin} is added at $a$.
\end{itemize}
Each of these merges adds one vertex for one fusion. The rule of Lemma~\ref{lem:merge} is the
standard graph rule for type-II fusion~\cite{Browne2005resource,Lee2023graph}; it was also checked
against both of our stabilizer engines on random inputs (Section~\ref{sec:method}).

\begin{proposition}[Universal lower bound]\label{prop:lb}
$\F(G)\ge n-3$ for every connected graph $G$ on $n\ge 3$ vertices.
\end{proposition}

\begin{proof}
Take a schedule with $g$ resources and $f$ fusions that produces $G$. Each resource has three
qubits and each fusion removes two, so $3g-2f=n$. Let $H$ be the multigraph whose vertices are
the resources and which has one edge for each fusion, joining the resources that hold its two
qubits (a loop if both are in one resource). Fusions in different components of $H$ act on
different tensor factors, so if two components contained surviving qubits the final state would
be a product across them, which is impossible since $|G\rangle$ is entangled across every
bipartition of a connected graph~\cite{Hein2006graphstates}. Hence one component, with $g_0$
resources and $f_0\ge g_0-1$ fusions, holds all survivors; every other component, with $g_i$
resources, has all its $3g_i$ qubits fused within it and so has $f_i=3g_i/2\ge g_i$ fusions.
Summing, $f\ge g-1$, and $n=3g-2f\le 3(f+1)-2f=f+3$.
\end{proof}

A counting form of the bound, in a merge model of schedules that also tracks fusions which split
or consume a component, is the Lean theorem \code{fusion\_cost\_lower\_bound}.
Since $3g-2f=n$, every schedule has $f\equiv n\pmod 3$, so $\F(G)\in\{n-3,\,n,\,n+3,\dots\}$.

\section{Proof of the characterization}\label{sec:proof}

The proof has a graph-theoretic core, which is a Lean theorem, and three steps that connect it to
qubits. Let $\Prod^{+}_k(G)$ say that $G$ is obtained from $P_3$ by $k$ steps, each adding a pendant
or a false twin at some vertex, interleaved with any number of local complementations and
relabelings (the Lean inductive \code{ProducibleByExt}). Let $\DHclass$ be the class of graphs obtained
from a connected three-vertex graph by pendant, false-twin and true-twin additions, closed under
isomorphism; these are exactly the connected distance-hereditary graphs on at least three
vertices~\cite{Bandelt1986distance} (the Lean inductive \code{PendantTwinBuildable}).

\begin{theorem}[Model characterization, Lean]\label{thm:model}
For a finite graph $G$, $\Prod^{+}_{|V(G)|-3}(G)$ holds if and only if $G\in\DHclass$
(\code{dh\_characterization}).
\end{theorem}

The forward implication is a structural induction on $\Prod^{+}$: pendant and false-twin
additions preserve $\DHclass$, and $\DHclass$ is closed under local complementation because rank-width at
most one is invariant under it~\cite{Oum2005rankwidth}. The converse follows the Bandelt-Mulder
construction of a distance-hereditary graph: a pendant addition is a pendant step, a false-twin
addition a false-twin step, and a true-twin addition is realized by local complementation, a
pendant step and local complementation, using the identity
\begin{equation}\label{eq:truetwin}
  \tau_u\big(\mathrm{addPendant}(G,u)\big)\;=\;\mathrm{addTrueTwin}\big(\tau_u(G),\,u\big)
\end{equation}
and the involutivity of $\tau_u$ (Section~\ref{sec:lean}).

To connect Theorem~\ref{thm:model} to schedules we need two lemmas, proved in
Appendix~\ref{app:normal}.

\begin{lemma}[Absorbing single-qubit Cliffords]\label{lem:frame}
Let $v$ be a vertex of a graph $G$ with at least one neighbor, and let $C$ be a single-qubit
Clifford unitary. Then there are a graph $G'$, obtained from $G$ by local complementations, and a
product $D$ of single-qubit Cliffords acting trivially on $v$, such that $C_v|G\rangle=D\,|G'\rangle$
up to a Pauli operator and a global phase.
\end{lemma}

\begin{lemma}[Single-component normal form]\label{lem:normal}
Suppose a schedule with $f$ fusions produces a connected graph $G$ on $n$ vertices. Then some
schedule with $f'\le f$ fusions and $g'$ resources, $3g'-2f'=n$, produces $G$ and has the following
form: one resource is the seed; each of the other $g'-1$ resources is attached by a single
\emph{merge}, a fusion between a qubit of the current component, which is a connected graph state,
and a qubit of the fresh resource with no single-qubit Clifford on either fused qubit; local
complementations of the component may occur between merges; after all merges come the remaining
$f'-g'+1$ fusions, each between two qubits of the component; single-qubit Cliffords on the
surviving qubits come last.
\end{lemma}

\begin{proof}[Proof of Theorem~\ref{thm:main}]
($\Leftarrow$) Let $G$ be connected and distance-hereditary on $n\ge 3$ vertices. Then $G$ is built
from $P_3$ or $K_3$, both $\LC$-equivalent to the resource, by $n-3$ pendant, false-twin and
true-twin additions~\cite{Bandelt1986distance}. By Lemma~\ref{lem:merge} each addition is one
merge with a fresh resource represented as $P_3$ (at a leaf or at the center) or as $K_3$. This
schedule has $n-3$ fusions, so $\F(G)\le n-3$, and Proposition~\ref{prop:lb} gives equality. The
graph-level part of this argument is the Lean theorem \code{dh\_forward}, which builds
$\Prod^{+}_{n-3}(G)$ from the Bandelt-Mulder construction, with true twins realized
through~\eqref{eq:truetwin}.

($\Rightarrow$) Suppose $\F(G)=n-3$ and apply Lemma~\ref{lem:normal} to an optimal schedule. Then
$f'\le n-3$, and Proposition~\ref{prop:lb} gives $f'=n-3$, so $g'=(n+2f')/3=n-2$ and the normal
form has $f'-g'+1=0$ fusions inside the component. The schedule is therefore a seed $P_3$ followed
by $n-3$ merges and local complementations. By Lemma~\ref{lem:merge} each merge adds a pendant,
a false twin or a true twin, and by~\eqref{eq:truetwin} a true-twin addition is a pendant step
between two local complementations. Hence the final graph $H$ satisfies $\Prod^{+}_{n-3}(H)$, and
$G\sim_{\LC}H$ because the closing single-qubit Cliffords do not change the $\LC$ class. By
Theorem~\ref{thm:model}, $H\in\DHclass$; distance-hereditary graphs are closed under local
complementation, so $G$ is distance-hereditary.
\end{proof}

\begin{remark}[An over-approximation in Lean]
The count used in the last proof, that a schedule at the bound has no fusion inside the
component, is also a Lean theorem in a model that is generous about such fusions. The inductive
\code{BlobSchedule} has the constructors of $\Prod^{+}$ and one more, \code{intra}, which costs
one fusion and may replace the current graph by \emph{any} graph with two fewer vertices.
Whatever a fusion of two qubits of one component does physically, it is an instance of this
constructor. The theorem \code{floor\_schedule\_iff\_dh} states that
$\mathsf{BlobSchedule}_{n-3}(G)$ holds if and only if $G\in\DHclass$: at the bound, a vertex count
(\code{BlobSchedule.natCard\_le}) excludes the \code{intra} step. Together with
Lemma~\ref{lem:normal}, which shows that every schedule is a \code{BlobSchedule} derivation up to
the free operations, this gives a second route to the reverse direction in which the only
unformalized step is the normal form itself.
\end{remark}

\section{The Lean development}\label{sec:lean}

The Lean development is in the directory \code{lean/EmpiricistLean} of the accompanying
repository, against \textsc{mathlib} at release \code{v4.31.0}. Table~\ref{tab:lean} in
Appendix~\ref{app:lean} lists the principal declarations and where they live. All of them use only
the axioms \code{propext}, \code{Classical.choice} and \code{Quot.sound}.

\paragraph{Production predicates.} Graphs are \textsc{mathlib} \code{SimpleGraph}s on an arbitrary
vertex type, and the production predicates carry exact graphs rather than equivalence classes.
The basic predicate $\Prod_k(G)$ has two constructors, the base $P_3$ and the \emph{leaf-merge}
that adds a pendant, closed under isomorphism; it has no local-complementation constructor and
generates exactly the trees. The extended predicate $\Prod^{+}_k$ adds the \emph{center-merge}, a
false-twin addition, and free local complementation. The leaf-merge is the definition
\code{ghz3LeafMerge}, and the lemma \code{ghz3CenterMerge\_iso\_addFalseTwin} identifies the
center-merge, written as the graph rewrite of Lemma~\ref{lem:merge}, with false-twin addition up
to isomorphism; the identity~\eqref{eq:truetwin} is \code{addPendant\_localComplement\_center}.
Inside Lean the merges are graph operations; their physical meaning is Lemma~\ref{lem:merge}.

\paragraph{Families.} Before the general theorem, the development proves $\F=n-3$ for natural
families, each as a statement that carries the general lower bound, quantified over arbitrary
schedules, alongside an explicit construction, so that no family theorem can be true for a
vacuous reason. Paths, stars and all trees are produced by leaf-removal induction; $K_n$ is one
local complementation from a star, and $K_{m,n}$ is three local complementations from a double
star. All of these graphs are distance-hereditary, so they are special cases of
Theorem~\ref{thm:main}.

\begin{remark}[Induction must track concrete graphs]\label{rem:concrete}
Although every statement about $\F$ is invariant under local complementation, the inductions
cannot be run on $\sim_{\LC}$ classes, because pendant addition does not descend to them:
$P_3\sim_{\LC}K_3$, yet a pendant at the center of $P_3$ gives the four-vertex star, while a
pendant at a vertex of $K_3$ gives the paw, which is $\LC$-equivalent to the path $P_4$ and not to
the star. The development therefore carries exact graphs through every induction and applies
local complementation only explicitly.
\end{remark}

\paragraph{Why the extended model.} $\Prod$ alone witnesses the bound only on the
local-complementation classes of trees, a strict subclass of the distance-hereditary graphs: of
the $175$ extremal orbits with $n\le 9$ (Section~\ref{sec:table}), $93$ contain a tree and $82$ do
not, two at $n=6$, four at $n=7$, nineteen at $n=8$ and fifty-seven at $n=9$. The missing
operation was the center-merge, a one-fusion primitive that the first model omitted. With it, the
forward direction is a structural induction on the Bandelt-Mulder construction that turns each
build step into one constructor, and the reverse direction is a structural induction on
$\Prod^{+}$.

\section{The certified tablebase}\label{sec:table}

\paragraph{Search.} Every connected graph with $3\le n\le 9$ was assigned to its
local-complementation orbit; there are $1, 2, 4, 11, 26, 101, 440$ orbits for $n=3,\dots,9$,
$585$ in all, in agreement with the published enumeration~\cite{Danielsen2006classification}.
The search follows Lemma~\ref{lem:normal}. Tier~0 is a breadth-first search over graphs in which
local complementation is free and a merge at any vertex, with a fresh resource fused at a leaf or
at the center, costs one fusion; true twins arise through~\eqref{eq:truetwin}. By
Lemma~\ref{lem:normal} it reaches exactly the orbits with $\F=n-3$. Tier~1 grows graphs by merges
to $n+2$ vertices and then performs one fusion of two qubits of the component; by
Lemma~\ref{lem:normal} the merges may come first, so this decides $\F=n$ for $n\le 7$. Every
fusion in the search is computed in the fixed basis of Section~\ref{sec:model}. For merges this
loses nothing by Lemma~\ref{lem:frame}; for the single fusion inside the component in Tier~1 we
re-ran the search with each of the six classes of single-qubit Clifford on one fused qubit
(which, together with the symmetry of the measured pair, covers all choices on both qubits; for
$n\le 5$ all $36$ pairs were tried) and no orbit changed. Each fusion outcome in a recorded
construction is computed by two stabilizer engines that share no code: a tableau simulator built on
\textsc{stim}~\cite{Gidney2021stim} and an independent implementation over GF(2). The resulting
table, $585$ rows each with an orbit representative, a verdict, a witness schedule when the value
is exact and a lower bound otherwise, is a single content-addressed artifact.

\begin{table}[t]
\centering
\begin{tabular}{lrrrrrrrr}
\toprule
$n$ & 3 & 4 & 5 & 6 & 7 & 8 & 9 & total\\
\midrule
connected orbits             & 1 & 2 & 4 & 11 & 26 & 101 & 440 & 585\\
$\F=n-3$ (extremal)          & 1 & 2 & 3 &  8 & 15 &  42 & 104 & 175\\
distance-hereditary          & 1 & 2 & 3 &  8 & 15 &  42 & 104 & 175\\
$\F=n$                       & 0 & 0 & 1 &  2 &  7 &   6 &   0 &  16\\
$\F=n+3$                     & 0 & 0 & 0 &  1 &  2 &   0 &   0 &   3\\
open, $\F\ge n$              & 0 & 0 & 0 &  0 &  0 &  53 & 336 & 389\\
open, $\F\ge n+3$            & 0 & 0 & 0 &  0 &  2 &   0 &   0 &   2\\
\bottomrule
\end{tabular}
\caption{Connected local-complementation orbits by number of vertices. The extremal orbits are
exactly the distance-hereditary ones. Exact values above the bound come from Tier~1 ($\F=n$ for
$n\le 7$) and from later certified constructions ($\F=n$ for six orbits with $n=8$, and the three
values $\F=n+3$, each with a lower bound from the exclusion of the lower tiers).}
\label{tab:counts}
\end{table}

\paragraph{Results.} Table~\ref{tab:counts} summarizes the table. A pendant-and-twin elimination
check~\cite{Bandelt1986distance}, run independently on all $585$ representatives, confirms that
the extremal orbits are exactly the distance-hereditary ones, with no exception; this is the
computational counterpart of Theorem~\ref{thm:main}. Nineteen orbits have exact values above the
bound. Sixteen are at $\F=n$. Three are at $\F=n+3$: a three-regular orbit on six vertices at
$\F=9$, realized with eight resources, nine fusions and three local complementations, and two
orbits on seven vertices at $\F=10$; their constructions were proposed by the model and accepted
by both engines, and the lower bounds come from Tier~0 and Tier~1.

\paragraph{Rank-width does not determine the cost above the bound.} Theorem~\ref{thm:main} says
that $\F=n-3$ exactly when the rank-width is one. A natural guess is $\F=n-3+3(\rw(G)-1)$. It
holds on all sixteen orbits with $\F=n$: none is distance-hereditary and each has linear rank-width
two, so each has rank-width exactly two. It fails at the next level. The six-vertex orbit with
$\F=9$ has rank-width two, because every graph on six vertices does (a branch decomposition whose
central node has three subtrees of two leaves each has width at most two) and the orbit is not
distance-hereditary; one of the two seven-vertex orbits with $\F=10$ has linear rank-width two,
hence rank-width two. For both the formula predicts $\F=n$, while the certified value is $n+3$. A
finer invariant is needed above the extremal class.

\section{Method: the harness and the discovery}\label{sec:method}

The results were obtained with Empiricist, an open-source harness that drives a language model
through research tasks under an explicit discipline. Its source, the Lean development, the
tablebase and the ledger of claims are in the accompanying repository.

\paragraph{A ledger of claims with evidence.} Every artifact (a computed value, a construction, a
conjecture, a Lean module) is recorded at a status in the ordered set
\textsc{refuted} $<$ \textsc{heuristic} $<$ \textsc{conjectured} $<$ \textsc{verified}$_n$ $<$
\textsc{certified} $<$ \textsc{formalized} and rises only with an evidence record from a certified
verifier; \textsc{refuted} is terminal, and nothing the model merely asserts is recorded above
\textsc{heuristic}. Artifacts are addressed by content hash and verifiers by the hash of their code
and of a golden test suite that they must pass before use.

\paragraph{The model never executes.} The model emits structured output, such as a schedule, a
conjecture or a complete Lean module, and the harness runs every check in a sandbox. This keeps
the boundary between a suggestion and a checked fact sharp.

\paragraph{Two stabilizer engines.} Fusion outcomes are computed by two engines that share no
code, and a construction is accepted only when both agree on the resulting state up to local
Clifford equivalence. The closed-form merge rule used in the hot loop of the search matches both
engines on randomized inputs ($400$ of $400$ for the center-merge), cross-checked against
canonical forms from \textsc{nauty}~\cite{McKay2014nauty}.

\paragraph{A hardened Lean gate.} A submitted Lean module is elaborated in a sandbox without
network or process creation, against a frozen read-only \textsc{mathlib}; the kernel result is
re-checked, the axioms are required to lie within $\{\code{propext},\code{Classical.choice},
\code{Quot.sound}\}$, every import must belong to an explicit trusted set, and the build
directory is swept after elaboration. The gate was developed against adversarial submissions
(poisoned object files, swaps between check and use, axiom-hiding elaborators) before it was used.

\paragraph{The discovery.} The harness gave the model the complete labeled catalogue of the $44$
orbits with $n\le 7$, each as an edge list tagged extremal or not, with no invariant, class name
or hint beyond the true statement that the target property is invariant under local
complementation. The model proposed that the extremal orbits are the distance-hereditary graphs,
gave the standard equivalent characterizations, and identified the smallest obstruction: the
unique non-extremal orbit on five vertices is the orbit of $C_5$. The proposal entered the ledger
as a conjecture, and a deterministic distance-hereditary check on all $585$ orbits, including the
$541$ orbits on eight and nine vertices that the model had not seen, agreed with it without
exception. The model then wrote the Lean proofs through a loop in which the gate returns the
remaining proof obligations and the model revises; it never ran Lean itself. The discipline also
caught two errors on the way: an arithmetic slip in an early hand-written note of the extremal
split ($185/400$ instead of $175/410$), and a family theorem whose first statement had the
construction class built into the claim, which the requirement to state the general lower bound
alongside each construction exposed.

\section{Related work}\label{sec:related}

\paragraph{Fusion cost in other resource models.} L\"obl \emph{et al.}~\cite{Lobl2025generating}
compute, by dynamic programming over local-complementation orbits, the minimum number of fusions
for all graph states on up to eight qubits when the inputs are caterpillar states generated by a
single quantum emitter, give a lookup table for about $2.8\times 10^{7}$ non-isomorphic graph
states, and derive bounds for larger states. Their model differs from ours in its resource, so
the values differ, and they do not characterize the graphs that attain a bound.
Lee and Jeong~\cite{Lee2023graph} optimize fusion-based generation of graph states with
graph-theoretical heuristics, and Ji \emph{et al.}~\cite{Ji2026motifs} decompose targets into
ring, star and linear motifs. We are not aware of a previous characterization of the graph states
that attain the lower bound in any of these models.

\paragraph{Graph states, local complementation and rank-width.} Local Clifford equivalence of
graph states is local complementation of graphs~\cite{VandenNest2004graphical,Bouchet1988transforming};
the orbits on up to twelve vertices were classified by Danielsen and
Parker~\cite{Danielsen2006classification} and mapped by Adcock
\emph{et al.}~\cite{Adcock2020mapping}. Rank-width~\cite{Oum2005rankwidth} is invariant under
local complementation, and distance-hereditary graphs, its first level, have appeared in quantum
information through the complexity of vertex-minor problems~\cite{Dahlberg2020transform} and the
classification of their local-equivalence classes~\cite{DH2026splitdecomp}.

\paragraph{Machine-assisted mathematics.} Language models have recently produced new
constructions in extremal combinatorics through program search~\cite{RomeraParedes2024funsearch,
Novikov2025alphaevolve}, and interactive theorem provers have been used to verify large
computer-assisted proofs~\cite{deMoura2021lean4,mathlib2020}. Here a model proposes a
characterization from data and writes the Lean proofs, and a harness admits only what a
verifier accepts.

\section{Discussion and open questions}\label{sec:discussion}

\paragraph{The boundary between Lean and the text.} The combinatorial optimization is proved in
Lean; the three qubit-level lemmas are proved in the text and corroborated by two stabilizer
engines and by the exhaustive search. A first-principles formalization would define graph states
and Pauli measurements in Lean and derive Lemma~\ref{lem:merge}; it would also formalize
schedules and the reordering of Lemma~\ref{lem:normal}. Both are larger projects than the result
itself.

\paragraph{Open questions.} First, a finer invariant governing $\F$ above the extremal class,
where rank-width is not enough (Section~\ref{sec:table}). Second, the complexity of computing
$\F$: Theorem~\ref{thm:main} makes the decision problem $\F(G)\stackrel{?}{=}n-3$ polynomial, since
distance-hereditary graphs are recognized in polynomial time by pendant-and-twin
elimination~\cite{Bandelt1986distance}, but the complexity of $\F$ in general is open. Third, the
analogous characterizations in other resource models, such as larger GHZ states, cluster-state
fragments or emitter-generated caterpillars.

\section*{Data and code availability}

The repository \url{https://github.com/sghowell/empiricist} (MIT license) contains the harness,
the Lean development, the tablebase and the claims ledger. The ancillary files of this paper
contain the Lean modules behind every declaration in Table~\ref{tab:lean} with a script that
builds them and prints their axioms, the tablebase, a script that re-derives every number in
Section~\ref{sec:table} using only the Python standard library, and the re-run of Tier~1 with
single-qubit Cliffords on the fused qubits, together with a README mapping each statement of the
paper to a file and a SHA-256 manifest.

\begin{acknowledgments}
This work was carried out with AI systems working under my direction. Inside the Empiricist
harness, the language model Fable~5 (Anthropic) classified the tablebase orbits, proposed the
distance-hereditary characterization and wrote the Lean proofs; it emitted only machine-checkable
artifacts and never executed code. Claude Code (Anthropic, model Claude Opus 5.5) revised the
manuscript, wrote the proofs of Lemmas~\ref{lem:merge}, \ref{lem:frame} and~\ref{lem:normal} and of
Proposition~\ref{prop:lb}, re-derived every number from the committed artifacts, and surveyed the
related literature. Codex (OpenAI) independently checked the proofs, the claims and the relation to
previous work; the errors it found have been corrected. I chose the problem and take full
responsibility for the content of this paper.
\end{acknowledgments}

\appendix

\section{Proofs of the qubit-level lemmas}\label{app:normal}

\begin{proof}[Proof of Lemma~\ref{lem:frame}]
Single-qubit Clifford unitaries modulo Pauli operators and phases form a group isomorphic to the
symmetric group on $\{X,Y,Z\}$, acting by conjugation. The unitaries $e^{-i\pi X/4}$ and
$e^{i\pi Z/4}$ act as the transpositions $(Y\,Z)$ and $(X\,Y)$, which generate it. By~\eqref{eq:lc},
local complementation at $v$ multiplies the local unitary on $v$ by $e^{-i\pi X_v/4}$, and local
complementation at a neighbor $u$ of $v$ multiplies it by $e^{i\pi Z_v/4}$. Neither operation
removes the edge $\{u,v\}$: $\tau_v$ changes only edges inside $\nbr(v)$ and $\tau_u$ only edges
inside $\nbr(u)$, and the edge $\{u,v\}$ is incident to the complemented vertex in both cases. So
$v$ keeps the neighbor $u$, and a suitable word in $\tau_v$ and $\tau_u$ gives a graph $G'$ and a
product $U=\bigotimes_w U_w$ of single-qubit Cliffords with $U|G\rangle=|G'\rangle$ up to a
Pauli operator and $U_v=C$ modulo Paulis. Then
$C_v|G\rangle=U_v|G\rangle=\big(\bigotimes_{w\neq v}U_w^{-1}\big)U|G\rangle$, which is
$D|G'\rangle$ with $D=\bigotimes_{w\ne v}U_w^{-1}$, up to a Pauli operator.
\end{proof}

\begin{proof}[Proof of Lemma~\ref{lem:normal}]
\emph{Step 1: single-qubit Cliffords.} A single-qubit Clifford on a qubit $q$ commutes with every
operation that does not act on $q$, and the only other operation acting on $q$ is the fusion that
consumes it, if any. Move each such gate to just after the preparation of its resource, to just
before the fusion that consumes $q$, or to the end of the schedule if $q$ survives. Gates moved to
the end do not change the $\LC$ class of the output. Gates on a fresh resource, together with the
choice of its graph, are absorbed below.

\emph{Step 2: discard what does not reach the output.} Let $H$ be the multigraph of the proof of
Proposition~\ref{prop:lb}. Components of $H$ without surviving qubits act on separate tensor
factors and can be deleted without changing the output, which does not increase the number of
fusions. As in that proof, all surviving qubits lie in one component, so after deletion $H$ is
connected, with $g'$ resources and $f'\le f$ fusions, and $3g'-2f'=n$.

\emph{Step 3: reorder.} Choose a spanning tree $T$ of $H$, a root resource, and a breadth-first
order of the resources from the root along $T$. For each non-root resource $c$ with parent $p$
in $T$, pick one fusion of $H$ joining $p$ and $c$; these are the merges. Reorder the fusions
so that the merges come first, in the breadth-first order of their child resources, followed by all
other fusions in their original relative order. Fusions act on disjoint pairs of qubits, since
each qubit is consumed at most once, so they commute as operations on the joint state: the final
state on the surviving qubits, conditioned on the outcomes, does not depend on the order, and a
different order changes only which Pauli corrections are applied. In the new order, when the merge
of $c$ is performed, its qubit in $p$ has not been consumed and $c$ has not yet been touched, so
the merge fuses a qubit of the component grown so far with a qubit of a fresh resource. Every
later fusion acts on two qubits of that component.

\emph{Step 4: remove the Cliffords on merge qubits.} Process the merges in order. Before the
merge of $c$, the component is the graph state of a connected graph $B$ on at least three vertices
up to single-qubit Cliffords on its qubits (the seed has three; each merge keeps the graph
connected by Lemma~\ref{lem:merge}, since both joined neighborhoods are non-empty, and adds one
vertex). Let $a$ be its fused qubit, carrying a Clifford $C$. Since $B$ is connected, $a$ has a
neighbor, and Lemma~\ref{lem:frame} replaces $C_a|B\rangle$ by $D|B'\rangle$ with $B'$ obtained by
local complementations and $D$ acting trivially on $a$; the factors of $D$ are carried forward on
the other qubits. On the fresh side, the resource with its gates is $E\,|P_3\rangle$ for a product
$E$ of single-qubit Cliffords; applying Lemma~\ref{lem:frame} to its fused qubit gives a graph
on three vertices, $P_3$ with some center or $K_3$, with no gate on the fused qubit. The merge is
now a fusion without single-qubit Cliffords on either fused qubit, and Lemma~\ref{lem:merge}
describes it. Gates carried onto qubits consumed by later merges are removed in the same way when
those merges are reached; gates on qubits consumed by the final fusions inside the component stay
where they are, as the statement allows; the rest move to the end.
\end{proof}

\section{Principal Lean declarations}\label{app:lean}

\begin{table}[h]
\centering
\begin{tabular}{lll}
\toprule
declaration & statement & file\\
\midrule
\code{fusion\_cost\_lower\_bound} & $f\ge n-3$ in a merge model of schedules & \code{Foundation.lean}\\
\code{tree\_min\_fusions} & $\F(T)=n-3$ for every tree $T$ & \code{TreeThm.lean}\\
\code{pathGraph\_min\_fusions} & paths & \code{Paper/PathGraph.lean}\\
\code{star\_min\_fusions} & stars & \code{Paper/Star.lean}\\
\code{complete\_min\_fusions} & $\F(K_n)=n-3$ & \code{Paper/Complete.lean}\\
\code{completeBipartite\_min\_fusions} & $\F(K_{m,n})=m+n-3$ & \code{Paper/CompleteBipartite.lean}\\
\code{ghz3CenterMerge\_iso\_addFalseTwin} & center-merge is false-twin addition & \code{CenterMerge.lean}\\
\code{addPendant\_localComplement\_center} & the identity~\eqref{eq:truetwin} & \code{TrueTwin.lean}\\
\code{dh\_forward}, \code{dh\_min\_fusions} & $\DHclass\Rightarrow\Prod^{+}_{n-3}$ & \code{DHCharacterization.lean}\\
\code{dh\_characterization} & $\Prod^{+}_{n-3}(G)\Leftrightarrow G\in\DHclass$ & \code{DHCharacterization.lean}\\
\code{floor\_schedule\_iff\_dh} & $\mathsf{BlobSchedule}_{n-3}(G)\Leftrightarrow G\in\DHclass$ & \code{Paper/FloorBridge.lean}\\
\bottomrule
\end{tabular}
\caption{Principal Lean declarations. Files named without a directory are library modules
(\code{lean/EmpiricistLean/EmpiricistLean/} in the repository); files in \code{Paper/} are
modules recorded as gate-certified artifacts in the claims ledger. The ancillary files contain
all of them as a self-contained Lake project with a script that builds it and prints the axioms of
every declaration.}
\label{tab:lean}
\end{table}

\bibliography{refs}

\end{document}
